\documentclass[11pt]{article}
\usepackage[margin=1in]{geometry}
\usepackage{amsmath,amssymb,amsthm,mathtools}
\usepackage{booktabs,longtable,tabularx,array}
\usepackage{enumitem}
\usepackage{microtype}
\usepackage{hyperref}
\usepackage{url}
\usepackage{xcolor}
\usepackage{fancyhdr}

\hypersetup{
  colorlinks=true,
  linkcolor=blue!55!black,
  citecolor=blue!55!black,
  urlcolor=blue!55!black,
  pdftitle={Independent Audit Note on Cobham-Free Full General Automatic Rigidity for Collatz-Invariant Sequences},
  pdfauthor={Muteki},
  pdfsubject={Post-release adversarial audit of the Full General ARC proof architecture},
  pdfkeywords={Collatz, automatic sequence, finite kernel, dyadic thinness, pumping, Lean 4, audit}
}

\pagestyle{fancy}
\fancyhf{}
\lhead{Full General ARC - Independent Audit Note}
\rhead{23 September 2026}
\cfoot{\thepage}
\setlength{\headheight}{15pt}

\newtheorem{proposition}{Proposition}[section]
\newtheorem{remark}[proposition]{Remark}
\newcommand{\N}{\mathbb{N}}
\newcommand{\Nzero}{\mathbb{N}_0}
\newcommand{\K}{\mathcal{K}}
\newcommand{\D}{\mathcal{D}}
\newcommand{\code}[1]{\path{#1}}

\title{\textbf{Independent Audit Note on\\
``Cobham-Free Full General Automatic Rigidity\\
for Collatz-Invariant Sequences''}\\[0.55em]
\large Adversarial examination of cross-base kernel transfer,\\
dyadic thinness, and topological pumping}

\author{Muteki\\\small Independent Researcher, Japan}
\date{\textbf{23 September 2026}\\\small Post-release supplementary audit note}

\begin{document}
\maketitle

\begin{abstract}
This note records a post-release adversarial audit of three structurally delicate
steps in the Cobham-free Full General ARC proof.  The audit was prompted by a
challenge that proposed possible failures in (i) the transfer from base-$2$
automaticity of a shortcut-Collatz-invariant sequence to base-$3$ automaticity
of its odd subsequence, (ii) the derivation of dyadic nowhere-thickness for the
odd-subsequence change set, and (iii) the theorem asserting that an odd-base
automatic Boolean support cannot be both infinite and dyadically nowhere thick.
A fourth concern was that the Lean theorem might hold only vacuously because
its hypotheses describe an empty class.

No counterexample to the published theorem is established in this note, and no
erratum or corrigendum is asserted.  Source-level tracing confirms that the
second step is implemented through an explicit affine pullback of a previously
proved local-density theorem, while the third step is not an unproved
cross-base heuristic: the formal development contains a dedicated chain from
odd-base pumping to dyadic residue permutation, child filling, and contradiction
with nowhere-thickness.  A proposed ternary Cantor-type support also fails to
supply the claimed counterexample in finite residue tests.  The vacuity concern
is ruled out constructively by explicit nonconstant-at-zero invariant automatic
sequences.  The first step is formally connected to the lower theorem
\code{arc_three_kernel_odd_finite}; however, because that lower proof predates
the present Cobham-removal package, a fresh line-by-line human re-derivation of
that lemma remains a sensible priority for continued audit.

Here ``independent'' means independent of the main packaged proof endpoint: this
is an additional adversarial review of the proof architecture, not an external
peer review by an independent third party.
\end{abstract}

\section{Status and purpose}

The parent paper proves the following rigidity statement for the shortcut
Collatz map
\[
T(n)=
\begin{cases}
 n/2, & n\equiv0\pmod 2,\\[2mm]
 (3n+1)/2, & n\equiv1\pmod 2.
\end{cases}
\]
For every base $B\ge2$, if $a:\Nzero\to\alpha$ has finite base-$B$ kernel and
satisfies
\[
a(T(n))=a(n)\qquad(n\in\Nzero),
\]
then $a$ is constant on the positive integers.  The value at zero may remain
independent.

This note does \emph{not} restate the complete proof.  Instead, it asks how the
argument could fail if the theorem were false.  The audit focuses on the three
places where two different structures meet:
\begin{enumerate}[label=(A\arabic*)]
\item base-$2$ finite-kernel structure versus a base-$3$ odd subsequence;
\item shortcut-Collatz dynamics versus $2$-adic cylinders;
\item odd-base finite automata versus dyadic nowhere-thickness.
\end{enumerate}
A fourth check asks whether the hypotheses are nonempty or whether the formal
theorem could be vacuous.

\begin{center}
\small
\renewcommand{\arraystretch}{1.15}
\begin{tabularx}{\textwidth}{@{}p{0.22\textwidth}X p{0.20\textwidth}@{}}
\toprule
\textbf{Audit target} & \textbf{Present finding} & \textbf{Status}\\
\midrule
A1. Odd subsequence becomes base-$3$ finite-kernel automatic
& The public Cobham-removal wrapper reduces exactly to
  \code{arc_three_kernel_odd_finite}, using only the even branch, odd branch,
  and finite base-$2$ kernel.  No Collatz-convergence hypothesis appears in
  the wrapper or in the recorded axiom audit.  A fresh standalone human proof
  of the lower lemma is still desirable.
& No confirmed defect; deeper audit retained.\\[0.7em]
A2. Change set is dyadically nowhere thick
& The source explicitly invokes the local-density theorem with gap $H=2$ on
  the affine image $x=2n+1$ and constructs a dyadic child in the original
  $n$-variable.
& Source-level route traced.\\[0.7em]
A3. Odd-base automatic + dyadically nowhere thick implies finite support
& The formal chain explicitly proves dyadic residue permutation, child filling,
  collision with nowhere-thickness, and then the finite-support theorem.
& Strongest remaining target for human re-derivation.\\[0.7em]
A4. Vacuous truth
& Explicit invariant finite-kernel sequences satisfying the hypotheses exist,
  including a family that may take a separate value at zero.
& Ruled out constructively.\\
\bottomrule
\end{tabularx}
\end{center}

\section{Audit A1: base-$3$ automaticity of the odd subsequence}

Let $a:\Nzero\to\alpha$ be shortcut-Collatz invariant and define
\[
b(n)=a(2n+1).
\]
The parent proof needs
\[
\K_2(a)\text{ finite}\quad\Longrightarrow\quad\K_3(b)\text{ finite}.
\]
The adversarial concern is that repeated interaction between coefficients $2$
and $3$ might generate infinitely many residual states, or might silently use
a global statement about Collatz orbits.

\subsection{What the active Cobham-removal wrapper actually assumes}

The active Stage-2 wrapper first derives the two branch identities
\begin{align}
a(2n)&=a(n),\\
a(2n+1)&=a(3n+2),
\end{align}
from shortcut invariance.  It then unfolds finite-kernel automaticity to obtain
\[
(\K_2(a)).\mathrm{Finite}
\]
and calls the lower theorem
\begin{center}\code{arc_three_kernel_odd_finite}.\end{center}
The output is precisely
\[
\left(\K_3\bigl(n\mapsto a(2n+1)\bigr)\right).\mathrm{Finite}.
\]
The resulting wrapper is
\begin{center}\code{arc2OddSubseq_threeAutomatic_of_twoAutomatic}.\end{center}

The recorded Lean axiom report for \code{arc_three_kernel_odd_finite} contains
only the standard logical/quotient dependencies
\begin{center}\code{propext}, \code{Classical.choice}, \code{Quot.sound}.\end{center}
Thus the active proof does not register a Collatz-convergence principle, a
Cobham principle, or an unproved custom axiom at this step.

\subsection{What this does and does not settle}

This source trace defeats one specific criticism: the Stage-2 proof is not
merely assuming that Collatz orbits ``eventually fall into finitely many
patterns.''  The finite base-$3$ kernel is supplied by a concrete lower Lean
theorem whose inputs are the branch identities and finite base-$2$ kernel.

However, formal dependency tracing is not the same as a new human proof.  The
lower theorem belongs to the earlier ARC2 kernel-transfer development and is
reused here.  Accordingly, the strongest responsible conclusion of this audit
is:
\begin{quote}
No defect has been identified in the active transfer step, but a fresh
line-by-line mathematical reconstruction of
\code{arc_three_kernel_odd_finite}, independent of its original development,
remains a high-value audit task.
\end{quote}
This is a retained audit item, not a correction notice.

\section{Audit A2: dyadic thinness of the odd-subsequence change set}

Define the adjacent-change set
\[
\D_b=\{n\in\Nzero:b(n)\ne b(n+1)\}.
\]
The concern is that shortcut Collatz dynamics involves multiplication by $3$,
whereas the thinness statement is expressed in dyadic residue classes.  A
purely heuristic appeal to ``compatibility'' between these structures would
indeed be insufficient.

\subsection{The actual affine pullback}

The active Stage-1 proof does not assume such compatibility.  It explicitly
applies the previously proved local-density theorem
\begin{center}\code{arc5_naturalInvariantBlock_in_dyadicCylinder}\end{center}
with gap parameter $H=2$ to the odd affine image.  Starting from a prescribed
parent class
\[
n\equiv r\pmod{2^m},
\]
the map
\[
x=2n+1
\]
sends it into
\[
x\equiv 2r+1\pmod{2^{m+1}}.
\]
Inside a deeper dyadic child, the local-density theorem supplies
\[
a(x)=a(x+2).
\]
After substituting $x=2n+1$, this is exactly
\[
a(2n+1)=a(2n+3),
\]
hence
\[
b(n)=b(n+1).
\]
The source then pulls the child cylinder back to the $n$-variable and proves
that it remains inside the original parent cylinder.  Therefore that child is
disjoint from $\D_b$.

The public endpoint is
\begin{center}\code{arc2OddSubseq_changeSet_nowhereThick},\end{center}
packaged as
\begin{center}\code{arc2CobhamRemoval_stage1_thinness}.\end{center}

\subsection{Audit conclusion for A2}

The objection ``base $2$ and multiplication by $3$ are transversal, therefore
the child need not exist'' does not identify a missing inference in the active
proof.  The existence of the child is precisely what the earlier local-density
machinery proves, and the affine conversion is carried out explicitly in the
Stage-1 source.

The natural next adversarial question is therefore not whether the two bases
are different, but whether the local-density theorem itself has been modeled
correctly and proved from the stated shortcut invariance.  That theorem was
already developed independently of automaticity in the ARC5 local-density
chain, so it can be audited as a separate arithmetic lemma family.

\section{Audit A3: odd-base automatic support versus dyadic nowhere-thickness}

This is the most structurally delicate point.  The relevant theorem states,
for odd $q\ge3$, that if a Boolean sequence has finite base-$q$ kernel and its
support is dyadically nowhere thick, then its support is finite.

At first sight this is surprising: the automaton reads base-$q$ digits, while
thinness is measured using powers of $2$.  The audit therefore asks whether the
proof actually bridges these incompatible scales, or merely assumes that it
can.

\subsection{The formal bridge is explicit}

The General ARC chain contains the following dedicated stages.
\begin{enumerate}[label=\arabic*.]
\item \textbf{Dyadic residue permutation.}  For an odd base, the normalized
      doubled-loop pumped family defines a map on residues modulo $2^M$.
      The theorem \code{arcGeneral_oddPumpResidueMap_injective} proves
      injectivity; finiteness of the common domain and codomain yields
      surjectivity; hence every residue modulo $2^M$ is hit.
\item \textbf{Dyadic child filling.}  The theorem
      \code{arcGeneral_evenLoopPumpedFamily_hits_every_child_of_exact_two_factor}
      lifts residue surjectivity into the assertion that the pumped family hits
      every deeper dyadic child of the relevant parent cylinder.
\item \textbf{Collision with nowhere-thickness.}  The theorem
      \code{arcGeneral_child_filling_not_nowhereThick} assumes a pumped family
      contained in the support, selects the child that thinness claims is empty,
      and then uses child filling to produce a support point in that same child,
      a contradiction.
\item \textbf{Infinite support gives a long pumpable word.}  Stage 9 extracts a
      sufficiently long canonical support word from any infinite automatic
      support and feeds it into the Stage-8 collision theorem.
\item \textbf{Finite-kernel handoff.}  Stage 10 converts finite kernel to an MSD
      DFAO and applies the preceding topological-pumping theorem.  Its public
      endpoint is
      \code{arcGeneral_kernelAutomatic_nowhereThick_support_finite}.
\end{enumerate}

Thus the base mismatch is not ignored; it is the central arithmetic object of
Stages 6--10.

\subsection{Adversarial Cantor-style test}

A natural proposed counterexample is the base-$3$ digit-restricted set
\[
C=\left\{\sum_{i\ge0}\varepsilon_i3^i:
\varepsilon_i\in\{0,1\},\text{ finitely many }\varepsilon_i\ne0\right\}.
\]
This set is base-$3$ automatic and infinite.  The proposed attack asks whether
it might also be dyadically nowhere thick.

A finite residue check gives the opposite indication.  Let
\[
C_L=\left\{\sum_{i=0}^{L-1}\varepsilon_i3^i:
\varepsilon_i\in\{0,1\}\right\}.
\]
For each modulus $2^m$ in the following range, a finite $L$ already makes
$C_L$ hit every residue class modulo $2^m$:

\begin{center}
\small
\begin{tabular}{@{}rrrrrrrrrr@{}}
\toprule
$m$ & 1 & 2 & 3 & 4 & 5 & 6 & 7 & 8 & 9 \\
\midrule
smallest tested $L$ & 1 & 3 & 4 & 5 & 7 & 8 & 9 & 11 & 12 \\
\bottomrule
\end{tabular}
\end{center}
For $m=10$, the smallest tested value is also $L=12$.  These are exact finite
computations of residue coverage, not random sampling.  They do not constitute
a proof for all $m$, but they show that this particular Cantor-style candidate
is not supported by the suggested finite test.  In particular, an earlier
finite search that reported failure to find an empty child in the parent
$0\pmod 2$ points against, rather than for, dyadic nowhere-thickness.

\subsection{Audit conclusion for A3}

No counterexample has been identified.  The formal proof contains an explicit
mixed-base mechanism whose purpose is to turn odd-base pumping into complete
coverage of dyadic children.  Nevertheless, because this mechanism is both
novel within the ARC development and mathematically nontrivial, A3 remains the
highest-priority target for a completely independent human proof written
without following the Lean file order.

\section{Audit A4: vacuous-truth concern}

The suggestion that the Lean theorem might be true only because no sequence
satisfies its hypotheses can be rejected constructively.

Fix arbitrary $x,y\in\alpha$ and define
\[
a(0)=x,\qquad a(n)=y\quad(n>0).
\]
Then $T(0)=0$, while $T(n)>0$ for every $n>0$.  Hence
\[
a(T(n))=a(n)
\]
for all $n\ge0$.  Moreover, in every base $B\ge2$, every kernel residual of
$a$ is either the original sequence $a$ or the constant sequence $n\mapsto y$.
Thus its base-$B$ kernel has at most two states.

If $x\ne y$, this sequence is not globally constant on $\Nzero$, but it is
constant on the positive integers exactly as the theorem predicts.  Therefore
the hypothesis class is manifestly nonempty, and the theorem is not a vacuous
implication caused by an inconsistent or empty premise.

\section{What a genuine counterexample would have to provide}

To refute the Full General ARC theorem, it is not enough to propose a coloring
that is constant on Collatz orbits, or a sparse automatic set, or a coloring
that respects only the even branch.  A genuine counterexample must explicitly
exhibit some $B\ge2$ and a sequence $a:\Nzero\to\alpha$ satisfying all three
conditions simultaneously:
\begin{enumerate}[label=(\roman*)]
\item the full base-$B$ kernel $\K_B(a)$ is finite;
\item $a(T(n))=a(n)$ for every $n\ge0$;
\item there exist positive $m,n$ with $a(m)\ne a(n)$.
\end{enumerate}
For example, assigning colors to hypothetical nonstandard Collatz components
might satisfy (ii), but it says nothing by itself about the finite-kernel
condition (i).  Likewise, extending a coloring by $a(2^kq)=a(q)$ enforces the
even relation but does not automatically enforce the odd relation
\[
a(2n+1)=a(3n+2).
\]
Any claimed counterexample must survive these exact checks.

\section{Lean 4: what formal verification does and does not certify}

The completed development reported successful project builds and no
\code{sorryAx} at the final audited endpoints.  The final theorem also no longer
carries an explicit \code{ARCCobhamPrinciple} argument or a
\code{[Finite alpha]} hypothesis.  This is meaningful evidence that the
packaged theorem follows from the definitions and imported lemmas accepted by
Lean.

Formal verification does not, however, replace the following external tasks:
\begin{itemize}
\item checking that the formal definitions express the intended mathematical
      notions;
\item checking the exposition of imported lower lemmas against an independent
      paper proof;
\item checking novelty and relation to the literature;
\item searching for an unintended theorem that is true for a different reason
      than the intended argument.
\end{itemize}
The present note is directed at the first two items.  Its role is therefore
supplementary: it records where the proof was attacked, what the source-level
trace showed, and which points still merit independent mathematical
reconstruction.

\section{Current audit judgment and next steps}

The adversarial review has not produced a confirmed error in the published
Full General ARC theorem.  It has instead isolated three high-value audit
points.  The present status is:
\begin{enumerate}[label=\arabic*.]
\item \textbf{A1: retained for deeper audit.}  Re-derive
      \code{arc_three_kernel_odd_finite} from its kernel definitions on paper,
      without invoking the packaged theorem.
\item \textbf{A2: source-level path confirmed.}  The affine dyadic pullback is
      explicit; a future audit should target the underlying local-density
      theorem itself rather than the mere fact that bases $2$ and $3$ differ.
\item \textbf{A3: highest priority.}  Reconstruct the residue-permutation and
      child-filling argument independently, and actively search for an infinite
      odd-base automatic support that could evade the proof.
\item \textbf{A4: closed.}  Vacuous truth is ruled out by explicit witnesses.
\end{enumerate}

Accordingly, this document is an \emph{Independent Audit Note}, not an erratum,
corrigendum, or replacement of the parent paper.  If a genuine defect is later
identified, it should be recorded separately and transparently as a correction.

\section*{Reproducibility note}

The finite Cantor-style residue experiment used the exact sets
\[
C_L\bmod 2^m
\]
formed by all $2^L$ choices of digits $\varepsilon_i\in\{0,1\}$, with set
compression after each added power of $3$.  The reported smallest $L$ values
are therefore exhaustive for the stated moduli.  They are included only as an
adversarial sanity check and are not used in the Lean proof.

\section*{Acknowledgment of AI assistance}

Generative AI systems were used to formulate adversarial objections, inspect
proof architecture, assist with Lean-source tracing, run finite computational
checks, and prepare this note.  An AI-generated objection is not treated as a
mathematical result unless its claim is independently verified.  Responsibility
for the released mathematical statements and the interpretation of the formal
artifacts remains with the human author.

\begin{thebibliography}{9}
\bibitem{Parent}
Muteki,
\emph{Cobham-Free Full General Automatic Rigidity for Collatz-Invariant Sequences},
Full General ARC release v1.0.0, 23 September 2026.

\bibitem{AlloucheShallit}
J.-P. Allouche and J. Shallit,
\emph{Automatic Sequences: Theory, Applications, Generalizations},
Cambridge University Press, 2003.

\bibitem{Cobham}
A. Cobham,
On the base-dependence of sets of numbers recognizable by finite automata,
\emph{Mathematical Systems Theory} \textbf{3} (1969), 186--192.
\end{thebibliography}

\end{document}